% SEGMENT OLP-0550-S01
\documentclass[../../../include/open-logic-section]{subfiles}

\begin{document}

\olfileid{sth}{ordinals}{wo}
\olsection{நன்கு வரிசைப்படுத்தல்கள்}

அடிப்படைக் கருத்து பின்வருவது:

\begin{defn}
பின்வரும் இரு நிபந்தனைகளையும் நிறைவுசெய்தால், அப்போது, அப்போது மட்டுமே,
$<$ என்ற தொடர்பு $A$ ஐ \emph{நன்கு வரிசைப்படுத்துகிறது}:
\begin{enumerate}
	\item $<$ ஒப்பிடத்தக்க பண்புடையது; அதாவது எல்லா $a, b \in A$
	க்கும் $a < b$ அல்லது $a = b$ அல்லது $b < a$;
	\item $A$ இன் ஒவ்வொரு வெற்றற்ற உட்கணத்திற்கும் ஒரு $<$-குறும
	!!{element} உண்டு; அதாவது $\emptyset \neq X \subseteq A$ எனில்
	$(\exists m \in X)(\forall z \in X)z \nless m$.
\end{enumerate}
\end{defn}

நாம் இப்போது பார்த்த மூன்று எடுத்துக்காட்டுகளும் உண்மையில் நன்கு
வரிசைப்படுத்தும் தொடர்புகளே என்பதை எளிதாகக் காணலாம்.

\begin{prob}
\Olref[sth][ordinals][idea]{sec} இயல் எண்கள்மீது மூன்று எடுத்துக்காட்டு
வரிசைப்படுத்தல்களை வழங்கியது. ஒவ்வொன்றும் ஒரு நன்கு வரிசைப்படுத்தல் என்று
சரிபார்க்கவும்.
\end{prob}

நன்கு வரிசைப்படுத்தலைப் பற்றிய சில தொடக்கநிலை, ஆனால் மிக முக்கியமான,
குறிப்புகள் இதோ.

\begin{prop}\ollabel{wo:strictorder}
$<$ ஆனது $A$ ஐ நன்கு வரிசைப்படுத்தினால், $A$ இன் ஒவ்வொரு வெற்றற்ற
உட்கணத்திற்கும் ஒரே ஒரு $<$-மீச்சிறு உறுப்பு உண்டு; மேலும் $<$ எங்கும்
தற்சுட்டற்றது, முழுச் சமச்சீரின்மை கொண்டது, கடப்புப் பண்புடையது.
\end{prop}

\begin{proof}
$X$ ஆனது $A$ இன் வெற்றற்ற உட்கணம் எனில், அதற்கு ஒரு $<$-குறும
!!{element} $m$ உண்டு; அதாவது $(\forall z \in X)z \nless m$.
$<$ ஒப்பிடத்தக்க பண்புடையதால், $(\forall z \in X)m \leq z$.
எனவே $m$ ஆனது $X$ இன் $<$-மீச்சிறு !!{element} ஆகும்.

எங்கும் தற்சுட்டின்மைக்காக, $a \in A$ ஐ நிலைப்படுத்துங்கள்;
$\{a\}$ இன் $<$-மீச்சிறு !!{element} $a$ என்பதால் $a \nless a$.
கடப்புப் பண்புக்காக, $a < b < c$ எனில், $\{a, b, c\}$ க்கு ஒரு
$<$-மீச்சிறு !!{element} இருப்பதால் $a < c$. எங்கும் தற்சுட்டின்மை,
கடப்புப் பண்பு ஆகியவற்றிலிருந்து முழுச் சமச்சீரின்மை பின்வருகிறது.
\end{proof}

\begin{prop}\ollabel{propwoinduction}
$<$ ஆனது $A$ ஐ நன்கு வரிசைப்படுத்தினால், எந்த வாய்பாடு $\phi(x)$
க்கும்:
\[
	\text{எனில் }(\forall a \in A)((\forall b < a)\phi(b) \lif 
		\phi(a))\text{, அப்போது }(\forall a \in A)\phi(a).
\]
\end{prop}

\begin{proof}
எதிர்மறைநேர்மாற்றை நிறுவுவோம். $\lnot(\forall a \in A)\phi(a)$,
அதாவது $X = \Setabs{x \in A}{\lnot\phi(x)} \neq \emptyset$ எனக்
கொள்வோம். அப்போது $X$ க்கு ஒரு $<$-குறும !!{element}, $a$, உண்டு.
எனவே $(\forall b < a)\phi(b)$, ஆனால் $\lnot \phi(a)$.
\end{proof}\noindent
இந்தக் கடைசிப் பண்பு, இயல் எண்கள்மீதான வலுவான தொகுத்தறிதல் கொள்கையை
உங்களுக்கு நினைவூட்ட வேண்டும்; அதாவது: $(\forall n \in
\omega)((\forall m < n)\phi(m) \lif \phi(n))$ எனில், $(\forall n \in
\omega)\phi(n)$. இந்தப் பண்பு நன்கு வரிசைப்படுத்தலை மிகவும்
\emph{வலுவான} கருத்தாக்குகிறது.\footnote{நினைவூட்டல்: வெளிப்படையாக
வேறுவிதமாகக் கூறப்படாதவரை, எல்லா வாய்பாடுகளிலும் அளவுருக்கள் இருக்கலாம்.}

\end{document}
